\subsection{Topological Origin of Fermionic and Bosonic Statistics}
  \label{app:relational_spin_statistics}

  The distinction between fermionic and bosonic behavior originates from
  the internal topological structure of~$\chi$ configurations in
  configuration space.

  \paragraph{Fermionic behavior.}
    Configurations with intrinsic $4\pi$-periodicity are double-valued under
    $2\pi$ reorientation.
    When projected, they exhibit fermion-like behavior: sign change under
    $2\pi$ rotations, restoration only after $4\pi$, and
    spin-$\tfrac{1}{2}$ transformation properties---without introducing
    fundamental spinors.

  \paragraph{Bosonic behavior.}
    Topologically orientable configurations return to an equivalent state
    after $2\pi$ reorientation, yielding integer-spin transformation
    properties.

  \paragraph{Topological mass ratios (heuristic).}
    For $\tilde{\chi}_c \approx 8.3$, the heuristic estimate
    $27\tilde{\chi}_c^2 \approx 1860$ agrees with the observed ratio
    $m_p/m_e = 1836.15$ at the $\sim\!1\%$ level.

    \medskip
    \noindent\textit{Status of this estimate.}
    No claim is made that the proton-to-electron mass ratio is predicted at the present
    stage (Appendix~B.8).
    The value $\tilde{\chi}_c \approx 8.3$ should be understood as an internally
    consistent, order-of-magnitude numerical estimate.
    Its role here is to provide a quantitative consistency check between a topological
    ansatz for composite sectors and the independently fixed normalization of the
    effective couplings.

    \medskip
    \noindent\textit{Heuristic ansatz and separation of assumptions.}
    The estimate rests on two distinct assumptions.

    \begin{enumerate}

      \item \textbf{Winding multiplicity factor.}
      Proton-like configurations are associated with a $w=3$ trefoil sector, while
      electron-like configurations correspond to $w=1$.
      A minimal multiplicity scaling yields a factor $w_p^3/w_e^3 = 3^3 = 27$.
      This factor is purely topological and does not depend on $\tilde{\chi}_c$.

      \item \textbf{Threshold-weight factor.}
      In addition to the winding multiplicity, we assume that the effective weight of an
      admissible sector carries a quadratic dependence on the dimensionless threshold
      parameter, schematically $\propto \tilde{\chi}_c^2$.
      This assumption is motivated by the appearance of $\chi_c$ in the mass scaling of
      localized configurations (Appendix~\ref{subsec:soliton_energy_mass},
      \ref{subsec:spectral-scaling-projection-ontology}--\ref{subsec:vacuum-energy-relaxation-capacity}), but is not
      derived here.

    \end{enumerate}

    \noindent
    Under these two assumptions one obtains the schematic estimate
    \begin{equation}
      \frac{m_p}{m_e}
      \;\approx\;
      \left(\frac{w_p}{w_e}\right)^3 \tilde{\chi}_c^2
      \;=\; 27\,\tilde{\chi}_c^2,
      \label{eq:massratio_topo}
    \end{equation}
    which should be viewed as a consistency check.
    A first-principles derivation requires (i) the tensorial Born--Infeld completion
    fixing the order-unity coefficient in Eq.~\eqref{eq:K0chic_gravity}, and (ii) a
    derivation of the sector weights from the admissible dynamics on
    $H^{(w)}_{\mathrm{adm}}$ as outlined in Section~\ref{sec:baryonic-sector}.

  \paragraph{Spin--statistics connection.}
    The relational-topological distinction between $4\pi$- and
    $2\pi$-periodic configurations provides a natural qualitative
    explanation of the spin--statistics connection, demonstrating that the
    observed dichotomy can arise from internal organization of~$\chi$ prior
    to any effective quantum description.
